\documentclass[11pt]{article}
\usepackage[utf8]{inputenc}
\usepackage[T2A,T1]{fontenc}
\usepackage{lmodern}
\usepackage[ukrainian,english]{babel}
\usepackage[margin=1in]{geometry}

\usepackage{amsmath,amssymb,amsfonts}
\usepackage{booktabs,multirow,array,tabularx}
\usepackage{graphicx,float}
\usepackage{tikz}
\usepackage{pgfplots}
\pgfplotsset{compat=1.18}
\usepackage[dvipsnames]{xcolor}
\usepackage{enumitem}
\usepackage{subcaption}
\usepackage{hyperref}
\hypersetup{
  colorlinks=true,
  linkcolor=RoyalBlue,
  citecolor=ForestGreen,
  urlcolor=BrickRed,
}
\usepackage[numbers,sort&compress]{natbib}

\title{Temporal Concept Drift in Legal Judgment Prediction:\\
Neural Baselines Across Three Epochs of Ukrainian Court Decisions}

\author{Volodymyr Ovcharov\\
\texttt{overthelex@legal.org.ua}}

\date{}

\begin{document}
\maketitle

% ============================================================
% ABSTRACT
% ============================================================
\begin{abstract}
Legal NLP benchmarks evaluate models on randomly split data, implicitly assuming that legal language is stationary. We test this assumption by fine-tuning four transformer encoders -- XLM-RoBERTa (base and large) and their legal-domain variants -- on Ukrainian court decisions from three temporal epochs defined by geopolitical disruptions: pre-war (2008--2013), hybrid war (2014--2021), and full-scale invasion (2022--2026). Each model is trained on one epoch and evaluated on all three, producing a $3 \times 3$ cross-temporal generalization matrix.

Three findings emerge. \textbf{(1)}~Forward degradation is severe: models trained on pre-war data lose \textbf{XX.X} percentage points of macro-F1 when applied to full-scale invasion era decisions, confirming and extending the 27.9-pp gap observed with classical baselines \citep{ovcharov2025tokenizer}. \textbf{(2)}~The degradation is asymmetric: backward transfer (full-scale$\rightarrow$pre-war) is substantially more robust than forward transfer, consistent with the hypothesis that legal language is additive -- new legal frameworks subsume older ones, but the reverse does not hold. \textbf{(3)}~Legal-domain pretraining (Legal-XLM-R) \textbf{does / does not} mitigate temporal drift compared to general-purpose XLM-R, suggesting that domain pretraining \textbf{does / does not} capture temporally stable legal representations.

We also evaluate continual learning as a mitigation strategy, finding that sequential fine-tuning across epochs \textbf{reduces / does not reduce} forward degradation compared to single-epoch training. These results establish the first neural temporal robustness benchmark for legal NLP and demonstrate that temporal drift is a dominant, underexplored source of performance degradation -- exceeding the impact of model selection and domain pretraining. The dataset (428K decisions across three epochs) is publicly available as a LEXTREME contribution.

\end{abstract}

% ============================================================
% 1. INTRODUCTION
% ============================================================
\section{Introduction}

The performance of NLP models degrades when the temporal distribution of test data diverges from training data -- a phenomenon known as temporal concept drift \citep{lazaridou2021temporal,luu2022temporal}. In general NLP, this manifests as outdated factual knowledge and shifting linguistic conventions. In the legal domain, the effect is amplified by three structural factors: legislative change introduces new statutes and modifies existing ones; judicial practice evolves as courts interpret new legislation; and exogenous shocks (reforms, conflicts) can alter the entire procedural framework within which courts operate.

Despite this, the dominant legal NLP benchmarks -- LexGLUE \citep{chalkidis2021lexglue}, LEXTREME \citep{niklaus2023lextreme}, and SCALE \citep{niklaus2023scale} -- evaluate models on randomly split data, treating temporal variation as noise rather than signal. This design choice is understandable for benchmarks focused on cross-lingual comparison, but it obscures a critical question for practitioners: \emph{how quickly does a deployed legal NLP model become unreliable?}

We address this question using a natural experiment. Ukraine's judicial system operated under three distinct regimes over the period 2008--2026:
\begin{enumerate}[nosep]
  \item \textbf{Pre-war (2008--2013):} Peacetime baseline. All 832 courts operational, stable procedural rules.
  \item \textbf{Hybrid war (2014--2021):} Crimea annexation (2014), loss of courts in occupied territories, judicial reform (2017), procedural modernization.
  \item \textbf{Full-scale invasion (2022--2026):} Martial law, altered procedural timelines, new Criminal Code articles (collaborationism, aiding the aggressor state), surge in military criminal cases.
\end{enumerate}

These three epochs are not arbitrary periodizations -- they correspond to structural breaks in citation network topology \citep{ovcharov2025citation}, 33--47\% decay in co-citation predictability \citep{ovcharov2025statute}, and 27.9-pp degradation in classical (TF-IDF) judgment prediction \citep{ovcharov2025tokenizer}. The present work extends these findings to neural models, answering the question: \emph{does fine-tuned transformer performance degrade in the same way, and can legal-domain pretraining or continual learning mitigate the effect?}

\paragraph{Contributions.}
\begin{enumerate}[nosep]
  \item We produce the first \textbf{neural cross-temporal generalization matrix} for legal judgment prediction, fine-tuning four XLM-R variants on 428K Ukrainian court decisions across three epochs.
  \item We quantify \textbf{forward--backward asymmetry} in neural temporal transfer and test whether legal-domain pretraining provides temporal robustness.
  \item We evaluate \textbf{continual learning} as a mitigation strategy for temporal drift in legal NLP.
  \item We release the \textbf{dataset} (428K decisions, three epochs, chronological splits) as a LEXTREME contribution -- the first Cyrillic-script subset with temporal annotations.
\end{enumerate}

% ============================================================
% 2. RELATED WORK
% ============================================================
\section{Related Work}

\paragraph{Legal NLP benchmarks.}
LexGLUE \citep{chalkidis2021lexglue} established a multi-task benchmark for English legal NLU, including case law from the European Court of Human Rights. LEXTREME \citep{niklaus2023lextreme} extended this to 11 datasets across 24 EU languages, evaluating XLM-R and domain-specific legal models. Both benchmarks use random train/test splits. SCALE \citep{niklaus2023scale} introduced longer-document tasks from the Swiss legal system, and FairLex \citep{chalkidis2022fairlex} added a fairness dimension. None of these benchmarks control for temporal distribution shift.

\paragraph{Legal judgment prediction.}
Swiss Judgment Prediction \citep{niklaus2021swiss} provided the first multilingual legal judgment prediction benchmark, with decisions spanning 2000--2020 but evaluated on random splits. PILOT \citep{cao2024pilot} introduced temporal pattern handling for case law retrieval, but focused on precedent identification rather than judgment prediction. Cross-X transfer \citep{niklaus2022cross} examined cross-lingual, cross-domain, and cross-regional transfer in legal NLP, but not cross-temporal transfer.

\paragraph{Legal-domain language models.}
LEGAL-BERT \citep{chalkidis2020legal} demonstrated the value of domain-specific pretraining for English legal text. The MultiLegalPile corpus \citep{niklaus2023multilegalpile} enabled pretraining of multilingual legal models (Legal-XLM-R), which achieved state-of-the-art on LEXTREME. SaulLM \citep{colombo2024saullm} scaled legal domain adaptation to 54B and 141B parameters. LeXFiles \citep{chalkidis2023lexfiles} provided a multinational English legal corpus with probing tasks. LEMUR \citep{ahmadi2026lemur} introduced multilingual legal embedding models for retrieval. A key question we address is whether legal-domain pretraining captures temporally stable representations.

\paragraph{Temporal generalization in NLP.}
\citet{lazaridou2021temporal} demonstrated that language models degrade on temporally shifted data, with performance inversely proportional to temporal distance. \citet{luu2022temporal} showed that temporal misalignment affects multiple NLP tasks beyond language modeling. \citet{dhingra2022time} proposed time-aware language models to mitigate temporal degradation. In the legal domain, \citet{ovcharov2025tokenizer} established a 27.9-pp forward degradation gap using TF-IDF classifiers on Ukrainian court decisions, but explicitly noted the absence of neural baselines as a limitation.

\paragraph{Continual learning.}
Catastrophic forgetting \citep{kirkpatrick2017overcoming} is a fundamental challenge when models are sequentially trained on new data distributions. We evaluate whether sequential fine-tuning across temporal epochs exhibits forgetting or knowledge accumulation.

% ============================================================
% 3. DATASET
% ============================================================
\section{Dataset}

\subsection{Source and Extraction}

We extract court decisions from the Unified State Register of Court Decisions (EDRSR, \foreignlanguage{ukrainian}{\textit{ЄДРСР}}), a publicly accessible database of all Ukrainian court decisions since 2006. The register contains over 100 million documents. We focus on civil and commercial jurisdictions, which provide the most consistent case structure across temporal epochs.

Each document is processed as follows: (1)~the facts section (\foreignlanguage{ukrainian}{\textit{встановив}}) is extracted as model input; (2)~the dispositive section (\foreignlanguage{ukrainian}{\textit{вирішив}}) is parsed for outcome classification; (3)~personally identifiable information is replaced with placeholder tokens (\texttt{[PERSON]}, \texttt{[ADDRESS]}, \texttt{[NUMBER]}). Texts are truncated to 10,000 characters.

\subsection{Temporal Epochs}

Decisions are divided into three epochs reflecting major geopolitical disruptions:
\begin{itemize}[nosep]
  \item \textbf{Pre-war (2008--2013):} 128,075 decisions. Peacetime judicial baseline.
  \item \textbf{Hybrid war (2014--2021):} 150,000 decisions. Post-Crimea annexation, judicial reform 2017.
  \item \textbf{Full-scale invasion (2022--2026):} 150,000 decisions. Martial law, procedural changes.
\end{itemize}

\subsection{Label Schema}

Outcomes are classified into three categories via regex extraction from the dispositive section:
\begin{itemize}[nosep]
  \item \textbf{Approved} (\foreignlanguage{ukrainian}{\textit{задоволено}}): claim fully granted.
  \item \textbf{Dismissed} (\foreignlanguage{ukrainian}{\textit{відмовлено}}): claim denied.
  \item \textbf{Partial} (\foreignlanguage{ukrainian}{\textit{частково задоволено}}): claim partially granted.
\end{itemize}

The pre-war epoch has a slight class imbalance (50K approved / 28K dismissed / 50K partial); the other two epochs are balanced at 50K per class. This imbalance reflects the natural distribution of outcomes in pre-war civil litigation.

\subsection{Chronological Splits}

Within each epoch, documents are split chronologically: the earliest 80\% form the training set, the next 10\% the validation set, and the most recent 10\% the test set. This prevents temporal leakage within epochs and ensures that models are always evaluated on decisions that postdate their training data.

\begin{table}[t]
\centering
\small
\caption{Dataset statistics. All splits are chronological within each epoch.}
\label{tab:dataset}
\begin{tabular}{lrrrr}
\toprule
\textbf{Epoch} & \textbf{Period} & \textbf{Train} & \textbf{Val} & \textbf{Test} \\
\midrule
Pre-war      & 2008--2013 & 102,460 & 12,807 & 12,808 \\
Hybrid war   & 2014--2021 & 120,000 & 15,000 & 15,000 \\
Full-scale   & 2022--2026 & 120,000 & 15,000 & 15,000 \\
\midrule
\textbf{Total} & 2008--2026 & 342,460 & 42,807 & 42,808 \\
\bottomrule
\end{tabular}
\end{table}

% ============================================================
% 4. EXPERIMENTAL SETUP
% ============================================================
\section{Experimental Setup}

\subsection{Models}

We evaluate four XLM-RoBERTa \citep{conneau2020xlmr} variants, covering the interaction of model scale and domain pretraining:

\begin{enumerate}[nosep]
  \item \textbf{XLM-R Base} (278M parameters) -- general multilingual baseline.
  \item \textbf{XLM-R Large} (560M parameters) -- scale comparison.
  \item \textbf{Legal-XLM-R Base} (278M) -- pretrained on the 689GB MultiLegalPile \citep{niklaus2023multilegalpile}.
  \item \textbf{Legal-XLM-R Large} (560M) -- legal-domain pretrained + scale.
\end{enumerate}

This $2 \times 2$ design (general vs.\ legal, base vs.\ large) isolates the effects of domain pretraining and model capacity on temporal robustness.

\subsection{Training Configuration}

All models are fine-tuned with the HuggingFace Transformers library \citep{wolf2020transformers} using the AdamW optimizer \citep{loshchilov2019decoupled}. Training configuration:

\begin{itemize}[nosep]
  \item Learning rate: $2 \times 10^{-5}$ (base), $1 \times 10^{-5}$ (large)
  \item Weight decay: 0.01
  \item Maximum sequence length: 512 tokens
  \item Training epochs: 5, with early stopping (patience 2) on validation macro-F1
  \item Warmup: 10\% of total steps
  \item Batch size: 16 per GPU (base), 8 per GPU (large), with gradient accumulation
  \item Hardware: 4$\times$ NVIDIA A10G GPUs (ml.g5.12xlarge)
\end{itemize}

Each experiment is run with three random seeds (42, 123, 456) and results are reported as mean $\pm$ standard deviation.

\subsection{Evaluation Protocol}

\paragraph{Experiment 1: In-epoch baselines.}
Each model is trained on epoch $E_i$ and evaluated on the test split of the same epoch. This establishes the diagonal of the generalization matrix: the best-case performance when train and test distributions match.

\paragraph{Experiment 2: Cross-epoch generalization.}
Each of the 12 trained models (4 models $\times$ 3 epochs) is evaluated on all three test splits, producing a $3 \times 3$ macro-F1 matrix per model. Key derived metrics:
\begin{itemize}[nosep]
  \item \textbf{Forward degradation:} $\Delta_{\text{fwd}} = F1(E_1 \rightarrow E_1) - F1(E_1 \rightarrow E_3)$
  \item \textbf{Backward degradation:} $\Delta_{\text{bwd}} = F1(E_3 \rightarrow E_3) - F1(E_3 \rightarrow E_1)$
  \item \textbf{Asymmetry gap:} $\Delta_{\text{fwd}} - \Delta_{\text{bwd}}$
\end{itemize}

\paragraph{Experiment 3: Continual learning.}
Models are sequentially fine-tuned on epochs in chronological order (forward: $E_1 \rightarrow E_2 \rightarrow E_3$) and reverse order (backward: $E_3 \rightarrow E_2 \rightarrow E_1$). After each stage, the model is evaluated on all three test sets. This tests whether sequential exposure mitigates or exacerbates temporal drift.

\paragraph{Experiment 4: Embedding drift analysis.}
We extract $[\text{CLS}]$ embeddings from each epoch-trained model on a shared sample of 1,000 documents per epoch. We compute pairwise cosine similarity and Centered Kernel Alignment (CKA; \citealt{kornblith2019similarity}) to quantify representational drift across temporal epochs.

\paragraph{Metrics.}
We report macro-F1 as the primary metric (robust to class imbalance), along with per-class F1 and accuracy. Statistical significance is assessed via paired bootstrap with 1,000 iterations.

% ============================================================
% 5. RESULTS
% ============================================================
\section{Results}

% --- PLACEHOLDER: Results will be filled after experiments ---

\subsection{In-Epoch Baselines}

\begin{table}[t]
\centering
\small
\caption{In-epoch performance (diagonal of generalization matrix). Macro-F1 (\%) averaged over 3 seeds.}
\label{tab:in_epoch}
\begin{tabular}{lcccc}
\toprule
\textbf{Model} & \textbf{Pre-war} & \textbf{Hybrid} & \textbf{Full-scale} & \textbf{Mean} \\
\midrule
XLM-R Base         & -- & -- & -- & -- \\
XLM-R Large        & -- & -- & -- & -- \\
Legal-XLM-R Base   & -- & -- & -- & -- \\
Legal-XLM-R Large  & -- & -- & -- & -- \\
\midrule
TF-IDF (baseline)  & 86.5 & 83.7 & 69.3 & 79.8 \\
\bottomrule
\end{tabular}
\end{table}

\textit{[Results to be filled after training runs complete.]}

\subsection{Cross-Epoch Generalization}

\begin{table}[t]
\centering
\small
\caption{Cross-epoch generalization matrix for XLM-R Base. Macro-F1 (\%) averaged over 3 seeds. Rows: training epoch; columns: test epoch. Bold: in-epoch (diagonal).}
\label{tab:cross_epoch_base}
\begin{tabular}{lccc}
\toprule
\textbf{Train $\backslash$ Test} & \textbf{Pre-war} & \textbf{Hybrid} & \textbf{Full-scale} \\
\midrule
Pre-war    & \textbf{--} & -- & -- \\
Hybrid war & -- & \textbf{--} & -- \\
Full-scale & -- & -- & \textbf{--} \\
\bottomrule
\end{tabular}
\end{table}

\textit{[One table per model, or a combined visualization. To be filled.]}

\subsection{Legal-Domain Pretraining Effect}

\textit{[Compare XLM-R vs Legal-XLM-R on temporal robustness metrics.]}

\subsection{Forward--Backward Asymmetry}

\textit{[Analysis of directional transfer asymmetry. Connection to ``legal language is additive'' hypothesis from \citet{ovcharov2025tokenizer}.]}

\subsection{Continual Learning}

\textit{[Sequential fine-tuning results. Catastrophic forgetting vs knowledge accumulation.]}

\subsection{Embedding Drift}

\textit{[CKA similarity, t-SNE visualizations. Connection to statute retrieval paper's 4.3\% semantic shift.]}

% ============================================================
% 6. DISCUSSION
% ============================================================
\section{Discussion}

\paragraph{Temporal drift exceeds model selection.}
\textit{[If confirmed: the forward degradation gap (XX pp) is larger than the difference between the best and worst model on any in-epoch evaluation -- consistent with the finding in \citet{ovcharov2025tokenizer} for classical and zero-shot baselines.]}

\paragraph{Legal-domain pretraining and temporal robustness.}
\textit{[Legal-XLM-R was pretrained on the MultiLegalPile, which spans multiple time periods. Does this temporal diversity confer robustness? Or does domain-specific knowledge become a liability when the legal framework changes?]}

\paragraph{Implications for benchmark design.}
Our results suggest that legal NLP benchmarks should report temporal robustness alongside aggregate performance. A model that achieves 90\% macro-F1 on a randomly split benchmark may be substantially less reliable when deployed on decisions from a later period. We propose that LEXTREME and future benchmarks include temporal split configurations alongside their standard evaluation.

\paragraph{Implications for practitioners.}
\textit{[Model maintenance cadence. Continual learning vs periodic retraining. Monitoring for temporal drift in production.]}

% ============================================================
% 7. LIMITATIONS
% ============================================================
\section{Limitations}

\paragraph{Single jurisdiction.}
Our findings are based on Ukrainian court decisions. While the three-epoch structure provides a natural experiment in temporal shift, the specific patterns of degradation may not generalize to other jurisdictions without comparable disruptions.

\paragraph{Single task.}
We evaluate judgment prediction (3-class classification). Other legal NLP tasks (NER, summarization, retrieval) may exhibit different temporal sensitivity profiles.

\paragraph{Encoder-only models.}
We focus on XLM-R variants, the standard LEXTREME evaluation architecture. Decoder-based models (GPT-4, Claude, Llama) and encoder-decoder architectures may show different temporal robustness properties.

\paragraph{Epoch boundaries.}
While our epoch boundaries correspond to documented structural breaks, the choice of exact cutoff dates (2014-01-01, 2022-01-01) is somewhat arbitrary. Sensitivity analysis with shifted boundaries would strengthen the findings.

\paragraph{Label noise.}
Outcome labels are extracted via regex from the dispositive section. While validated on a 273-document subset \citep{ovcharov2025tokenizer}, systematic errors in the regex parser could introduce epoch-dependent label noise.

% ============================================================
% 8. CONCLUSION
% ============================================================
\section{Conclusion}

We have presented the first neural temporal robustness benchmark for legal NLP, fine-tuning four XLM-R variants on 428K Ukrainian court decisions across three temporal epochs. Our cross-epoch generalization matrices reveal that temporal drift is a dominant source of performance degradation, exceeding the impact of model selection and domain pretraining.

The forward--backward asymmetry in temporal transfer provides mechanistic insight: legal language is additive, with newer frameworks subsuming older ones but not vice versa. This has direct implications for model maintenance -- backward-compatible models trained on recent data are more temporally robust than forward-deployed models trained on historical data.

We release the full dataset (428K decisions, three epochs, chronological splits) as a LEXTREME contribution, providing the first Cyrillic-script legal NLP subset with temporal annotations. We advocate for temporal robustness evaluation as a standard component of legal NLP benchmarking.

\paragraph{Data and code availability.}
The dataset is available at \url{https://huggingface.co/datasets/overthelex/ukrainian-court-decisions}. Training scripts and evaluation code are released at \url{https://github.com/overthelex/SecondLayer}.

% ============================================================
% REFERENCES
% ============================================================
\bibliographystyle{plainnat}
\bibliography{references}

\end{document}
